\textbf{(CAP-COLUNI 2024 -- Q13)} Em 19 de janeiro de 2006, a sonda da missão New Horizons da NASA partiu de Cabo Canaveral, na Flórida-USA, em direção ao planeta Plutão. O objetivo da missão era estudar a geologia do planeta, as luas e objetos do Cinturão de Kuiper, nos limites do Sistema Solar. A sonda chegou ao planeta destino, em meados de 2015. Foi uma viagem de quase 5 bilhões de quilômetros, que durou exatos 9 anos e meio. De acordo com a agência espacial, a viagem ``levou cerca de um minuto a menos do que o previsto''.

HOUSEL, Morgan. A psicologia financeira. 1. Ed. Rio de Janeiro. Harper Collins, 2021. p. 250.\\
Disponível em: www.agenciabrasil.com.br. Acesso em: 29 ago 2023 (adaptado).

O cálculo da Velocidade Média (VM) de um objeto é obtido por meio da razão entre a distância que ele percorre e o tempo gasto nesse percurso. Considere um ano com 360 dias e um mês com 30 dias. A alternativa que mais se aproxima da VM da sonda New Horizons, na viagem de Cabo Canaveral até Plutão, é:

\begin{enumerate}
	\item 62.000 km/h.
	\item 16.000 km/h.
	\item 31.000 km/h.
	\item 8.000 km/h.
\end{enumerate}
